\begin{figure}[!t]
\centering
\includegraphics[width=\textwidth]{imgs/overview.pdf}
\caption{\textbf{Overview of the controllable image-to-video generation framework.}
To control trajectories of scene elements, we optimize the latent $\bm{z}_t$ at specific denoising timesteps $t$ of a pre-trained video diffusion model.
First, we extract semantically aligned feature maps from the denoising U-Net to estimate the video layout.
Next, we enforce cross-frame feature similarity along the bounding box trajectory to drive the motion of each region.
To preserve the visual quality of the generated video, a frequency-based post-processing method is applied to retain high-frequency noise of the original latent $\bm{z}_t$. The updated latent $\tilde{\bm{z}}_t$
 is input to the next denoising step.
}
\label{fig:method}
\vspace{-2.0mm}
\end{figure}


\section{Method}
\label{sec:method}

In this section, we describe our method for the trajectory control task in image-to-video generation (\cref{sec:3_task}).
Our framework, SG-I2V, builds on the publicly available image-to-video diffusion model Stable Video Diffusion~(SVD)~\citep{blattmann2023stable},
and consists of two main steps. 
First, we extract and semantically align the feature maps from a specific layer of SVD during the early steps of the diffusion process (\cref{sec:motion_estimation}); we show that such feature maps are especially effective at influencing motion in the output video.
Second, we optimize the noisy latent (i.e., the input to the denoising network) to enforce similarity between features within the bounding box trajectories (\cref{sec:motion_optimization}).
However, we find that naive optimization of latent is prone to overfitting and often results in low-quality generation.
Thus, we employ frequency-based post-processing to retain an in-distribution noisy latent (\cref{sec:fft}).
Our entire pipeline is summarized in \cref{fig:method}.

\subsection{Trajectory Control in Image-to-Video Generation} %
\label{sec:3_task}
The goal of our work is to build a zero-shot framework that takes an input image $\mathbf{I} \in \mathbb{R}^{H \times W \times 3}$ (with height $H$, width $W$) and generates a video with $N$ frames, where elements in the input image move in a user-specified fashion.
Inspired by DragAnything~\citep{wu2024draganything}, we further assume that a user provides a set of $B$ input bounding box trajectories $\{\mathcal{B}_b\}_{b=1}^B$, each parameterized by a height $h_b$ and width $w_b$, as well as center point coordinates  for each output video frame: $\mathbf{c}_{b, n}\in\mathbb{R}^2, 1\leq n \leq N$.
For simplicity, we assume bounding boxes cannot extend outside the image, and we denote the $b$-th bounding box in the $n$-th frame as $\mathcal{B}_{b, n} = \{h_b, w_b, \mathbf{c}_{b,n}\}$. 

Our aim is to constrain regions of the input image falling within a bounding box to follow the trajectory of the same bounding box in the output video.
Thus, the motion of dynamic foreground objects can be controlled by placing bounding boxes around them and specifying the desired trajectory.
On the other hand, camera motion can be specified by placing bounding boxes on static background regions and specifying a trajectory opposite to the desired camera movement. %
Further, we can also set the bounding box trajectory to the zero vector to keep the region static.
Overall, this formulation provides intuitive and unified control over object and camera motion, which are sometimes treated separately in previous controllable image-to-video frameworks~\citep{MotionCtrl, DirectAVideo, li2024image}.





\subsection{Extracting Semantic Video Layout}
\label{sec:motion_estimation}
\paragraph{Preliminaries: Stable Video Diffusion.}
Video diffusion models~\citep{VideoDiffusionModels} learn a data distribution $\bm{x}_0 \sim p_\theta(\bm{x}_0)$ by gradually denoising a video corrupted by Gaussian noise. 
The output denoised video is thus drawn from the distribution $p_\theta(\bm{x}_0) = \int p_\theta(\bm{x}_{0:T}) \ d\bm{x}_{1:T}$, where $\bm{x}_0 \in \mathbb{R}^{N \times H \times W}$ is a clean video, and $\bm{x}_{1:T}$ are intermediate noisy samples.
For simplicity, we omit the channel dimension throughout the paper.
To reduce computation, Stable Video Diffusion (SVD)~\citep{blattmann2023stable} performs the diffusion process in a latent space, where a variational autoencoder~\citep{VAE} maps a raw video $\bm{x}_0$ to a latent $\bm{z}_0 \in \mathbb{R}^{N \times h \times w}$. \\
Since this work aims to animate an existing image, we utilize the image-to-video variant of SVD, which concatenates a conditioning frame with noisy latent ($\bm{z}_t$) and runs a 3D U-Net~\citep{U-net} to predict the noise.
The 3D U-Net contains a downsampling and an upsampling path.
Specifically, the upsampling path consists of three stages operating at different resolutions, where each stage contains three blocks with interleaved residual blocks~\citep{ResNet}, spatial, and temporal attention layers~\citep{Transformer}.
We will call these three stages bottom, middle, and top from lower to higher resolution.
For more details, we refer readers to the original paper of SVD~\citep{blattmann2023stable}.


\vspace{-2.0mm}
\paragraph{SVD feature map analysis.} 
In image diffusion models, prior work has shown that output feature maps of upsampling blocks in the middle stage of the denoising U-Net are \textit{semantically aligned}~\citep{tang2023emergent, namekata2024emerdiff, hedlin2024unsupervised, zhang2023tale, luo2023dhf}, i.e., regions belonging to the same object tend to have similar feature vectors. %
Such semantically aligned feature maps are useful in estimating the layout of generated images, enabling spatial control of objects~\citep{shi2024dragdiffusion, mou2023dragondiffusion}.
Therefore, we first examine whether SVD feature maps are also semantically correlated across \textit{both spatial and temporal dimension}. %
\cref{fig:correspondence} visualizes the principal components of feature maps extracted from the upsampling block and spatial attention layers.
We observe that SVD feature maps exhibit weak semantic correspondence across frames at early denoising steps, leading to inaccurate object trajectory estimation.
Yet, we want to operate at early steps as they decide the structure of generated videos~\citep{materzynska2023customizing}.
This dilemma prompts us to align these features before applying optimization.

\textbf{Feature alignment with modified self-attention.} 
SVD leverages separate spatial and temporal self-attention to model the entire video.
Since spatial self-attention is only applied per frame, it does not produce cross-frame aligned features.
While temporal attention communicates across frames, it only attends to the same pixel position on the feature map, which may be inadequate for capturing semantic information spatially. %
To address this issue, inspired by \citep{Tune-a-video}, we modify the spatial self-attention on each frame to directly attend to the first frame.
Concretely, for the $n$-th frame, the original spatial self-attention works as $\bm{F}_n = 
\mathrm{Softmax}(\frac{\bm{Q}_{n} \cdot \bm{K}_{n}^T}{\sqrt{D}})\cdot \bm{V}_{n}$, where $\bm{F}_n$ %
is the %
outputs of self-attention, $\bm{Q}, \bm{K}, \bm{V}$ are the query, key, and value tokens, respectively, and $D$ is the dimensionality of the key and query tokens~\cite{Transformer}.
Instead, we replace the key $\bm{K}_{n}$ and value $\bm{V}_{n}$ of each frame with $\bm{K}_{1}$ and $\bm{V}_{1}$ from the first frame, leading to a new operation $\tilde{\bm{F}}_n = 
\mathrm{SoftMax}(\frac{\bm{Q}_{n} \cdot \mathrm{SG}(\bm{K}_{1})^T}{\sqrt{D}}) \cdot \mathrm{SG}(\bm{V}_{1})$.
We apply a stop gradient $\mathrm{SG}(\cdot)$ on $\bm{K}_{1}$ and $\bm{V}_{1}$ to stabilize the subsequent optimization process.
Now, all the modified %
feature maps $\tilde{\bm{F}}_n$ are weighted combinations of $\bm{V}_1$, exhibiting a stronger cross-frame correspondence while still maintaining the object layout of each frame, as shown in the bottom row of \cref{fig:correspondence}. Notably, this modification occurs during the loss computation only and does not affect the denoising steps.


\subsection{Trajectory Control with Latent Optimization}
\label{sec:motion_optimization}


So far, we have obtained spatio-temporally aligned feature maps $\tilde{\bm{F}}_n(z_t) \in \mathbb{R}^{h \times w \times d}$ at each frame given the noisy latent $\bm{z}_t$ as input (for simplicity, we resize $\tilde{\bm{F}}_n$ to the same resolution as the noisy latent).
Recall that our goal is to control the output video frames so that the bounding boxes $\mathcal{B}_b$ identified in the first frame move along the associated trajectories.  
Inspired by prior drag-based control methods~\citep{shi2024dragdiffusion,pan2023drag}, we optimize the noisy latent $\bm{z}_t$ to enforce cross-frame similarity between features within bounding boxes. 
The optimization objective is as follows:
\begin{equation}\label{eq:feature-optimization}
    \bm{z}_t^* = \argmin_{\bm{z}_t} \sum_{b\in[1, B], n\in[2, N]}\lVert \bm{G}_{b} \odot \left(\tilde{\bm{F}}_n(\bm{z}_t)[\mathcal{B}_{b, n}] - \mathrm{SG}(\tilde{\bm{F}}_1(\bm{z}_t)[\mathcal{B}_{b, 1}])\right) \rVert_2,
\end{equation}
where $\odot$ is Hadamard product, and $\tilde{\bm{F}}_n(\bm{z}_t)[\mathcal{B}_{b, n}] \in \mathbb{R}^{h_b \times w_b \times d}$ is feature maps cropped by the bounding box $\mathcal{B}_{b, n}$. 
Following DragAnything~\citep{wu2024draganything}, we weight the feature difference using a Gaussian heatmap $\bm{G}_{b} \in \mathbb{R}^{h \times w}$. This focuses on optimizing pixels closer to the bounding box center, as pixels near the edge may be background pixels that we do not want to move. 
\vspace{-2.5mm}
\paragraph{Selective latent optimization.}
We optimize \cref{eq:feature-optimization} on a subset of denoising timesteps and self-attention layers.
Concretely, we only select early denoising timesteps as the coarse structure of output frames is determined at these timesteps~\citep{DMProcessAnalyze,materzynska2023customizing}.
In addition, consistent with previous works in image diffusion models~\citep{shi2024dragdiffusion,mou2023dragondiffusion}, we observe that feature maps extracted from the middle stage of the denoising U-Net are more semantically correlated, and thus we use them for optimization.

\vspace{-2.0mm}
\subsection{High-Frequency Preserved Post-Processing}
\label{sec:fft}

Although the presented pipeline already enables trajectory control in video generation, we notice a quality degradation in generated videos.
We attribute this degradation to the deviation of $\bm{z}_t^{*}$ from the sampling distribution of the diffusion process after optimization. 
A recent work FreeInit~\citep{wu2023freeinit} observed that motions of generated videos are mostly encoded in the low-frequency component of noisy latent.
Inspired by this, we propose to discard the high-frequency component of the optimized latent $\bm{z}_t^*$ and replace it with the high-frequency component of the original latent $\bm{z}_t$. 
Formally, we obtain the new latent $\tilde{\bm{z}}_t$ as follows:
\begin{equation}
\begin{aligned}
    \tilde{\bm{z}}_t = \textsc{IFFT}_{\mathrm{2D}}\left(\textsc{FFT}_{\mathrm{2D}}(\bm{z}_t^*) \odot \mathbf{H}_\gamma + \textsc{FFT}_{\mathrm{2D}}(\bm{z}_t) \odot (1 - \mathbf{H}_\gamma)\right),
\end{aligned}
\end{equation}
where $\textsc{FFT}_{\mathrm{2D}}$ is the Fast Fourier transformation applied to each frame, and $\textsc{IFFT}_{\mathrm{2D}}$ is the corresponding inverse operation.
We set $\mathbf{H}_\gamma$ as the frequency response of a 2D low-pass filter (we follow FreeTraj and use a Butterworth filter) with cut-off frequency $\gamma$.
This post-processing step retains the target motion signals encoded in the low-frequency component of $\bm{z}^*_t$ while eliminating undesirable high-frequency disruptions. %
